\documentclass[aspectratio=169,10pt,english]{beamer}

\input{../../../_preambulo_beamer.tex}
\input{../../../_comandos_beamer.tex}

\selectlanguage{english}

\title{MA1001B - Statistical Modeling for Decision Making}
\subtitle{Lesson 28: Student t Interval when sigma is Unknown}
\author{}
\institute{Tecnol\'ogico de Monterrey}
\date{}

\begin{document}

\begin{frame}[plain]
  \vspace{1.2cm}
  {\Large\bfseries MA1001B - Statistical Modeling for Decision Making\par}
  \vspace{0.7cm}
  {\large Lesson 28: Student t Interval when sigma is Unknown\par}
  \vspace{0.7cm}
  {\color{TecRojo}\rule{\textwidth}{1.2pt}}
  \vspace{0.8cm}

  MA1001B Faculty\\[0.3cm]
  Term\\[0.3cm]
  Tecnol\'ogico de Monterrey
\end{frame}

\begin{frame}{The Unknown-Sigma Interval Question}
  \begin{alertblock}{Guiding question}
    How is a 95\% interval for an unknown mean built when $\sigma$ must be
    replaced by the sample $s$?
  \end{alertblock}

  \vspace{0.6em}
  By the end of this lesson, you can:
  \begin{itemize}
    \item compute $s$ and $\widehat{\mathrm{SE}}=s/\sqrt{n}$;
    \item obtain $t^{*}$ from the $t$ distribution with $\mathrm{df}=n-1$;
    \item form $\bar{x}\pm t^{*}\cdot s/\sqrt{n}$;
    \item explain how one outlier inflates $s$ and the interval.
  \end{itemize}

  \vfill
  \scriptsize All fill volumes used today are synthetic. No real company data are used.
\end{frame}

\begin{frame}{The Same Sample, Now with Unknown $\sigma$}
  \small
  \begin{center}
    \begin{tabular}{lr}
      \toprule
      \textbf{Quantity} & \textbf{Value} \\
      \midrule
      Sample size $n$ & 36 bottles \\
      Degrees of freedom $n-1$ & 35 \\
      Point estimate $\bar{x}$ (seed 42) & 502.738498 ml \\
      Sample $s$ & 8.390375 ml \\
      \bottomrule
    \end{tabular}
  \end{center}

  \vspace{0.5em}
  \begin{exampleblock}{Estimated standard error}
    \[
      \widehat{\mathrm{SE}}=\frac{s}{\sqrt{n}}=\frac{8.390375}{\sqrt{36}}=1.398396.
    \]
    The extra uncertainty in $s$ is why the critical value is $t$, not $z$.
  \end{exampleblock}
\end{frame}

\begin{frame}[fragile]{Step 1: Sample Standard Deviation}
  \begin{columns}[T]
    \begin{column}{0.42\textwidth}
      \small
      \begin{lstlisting}
s = np.std(sample, ddof=1)
se_hat = s / np.sqrt(n)
df = n - 1
      \end{lstlisting}

      \begin{block}{Verified output}
        $\mathrm{df}=35$\\
        $s=8.390375$ ml\\
        $\widehat{\mathrm{SE}}=1.398396$
      \end{block}
    \end{column}
    \begin{column}{0.58\textwidth}
      \centering
      \includegraphics[width=\textwidth]{../figures/t_ci_mean_01_sample_sd.png}
      \vspace{0.2em}
      \scriptsize Histogram and estimated SE band from Step 1
    \end{column}
  \end{columns}
\end{frame}

\begin{frame}{Replace $z^{*}$ by $t^{*}$ with $\mathrm{df}=35$}
  \small
  From \texttt{scipy.stats.t.ppf}:
  \[
    t^{*}=t_{0.975,35}=2.030108 > z^{*}=1.96.
  \]
  Margin of error and interval:
  \[
    E=2.030108\times 1.398396=2.838894,
  \]
  \[
    \bar{x}\pm E=[499.899604, 505.577392].
  \]
  This seed-42 interval covers the hidden $\mu=502$. The $t$ critical value
  is larger than $1.96$; $s$ itself may still be smaller than $\sigma$.
\end{frame}

\begin{frame}[fragile]{Step 2: The 95\% t Interval}
  \begin{columns}[T]
    \begin{column}{0.40\textwidth}
      \small
      \begin{lstlisting}
t_star = stats.t.ppf(
    0.975, df=35
)
margin = t_star * se_hat
      \end{lstlisting}

      \begin{block}{Verified output}
        $t^{*}=2.030108$\\
        Margin $=2.838894$\\
        CI $[499.899604, 505.577392]$
      \end{block}
    \end{column}
    \begin{column}{0.60\textwidth}
      \centering
      \includegraphics[width=\textwidth]{../figures/t_ci_mean_02_t_interval.png}
      \vspace{0.2em}
      \scriptsize Unknown-sigma t interval from Step 2
    \end{column}
  \end{columns}
\end{frame}

\begin{frame}{Step 3: One Outlier Inflates $s$ and the Interval}
  \begin{columns}[T]
    \begin{column}{0.42\textwidth}
      \small
      One fill replaced by $560$ ml.

      \vspace{0.4em}
      $s$ rises from $8.390375$ to $12.616189$.

      \vspace{0.4em}
      Contaminated interval:
      \[
        [499.767301, 508.304710].
      \]
      Width grows by $2.859620$ ml.

      \vspace{0.5em}
      \textbf{Limit:} one extreme observation inflates $s$ and the $t$
      interval. The formula still runs; the scale estimate does not.
    \end{column}
    \begin{column}{0.58\textwidth}
      \centering
      \includegraphics[width=\textwidth]{../figures/t_ci_mean_03_outlier_inflates_s.png}
      \vspace{0.2em}
      \scriptsize Clean versus outlier-inflated t intervals from Step 3
    \end{column}
  \end{columns}
\end{frame}

\begin{frame}{Concept Check}
  \begin{enumerate}
    \item Why is the critical value $t_{0.975,35}$ instead of $1.96$?
    \item Compute $\widehat{\mathrm{SE}}=8.390375/\sqrt{36}$.
    \item Why can this $t$ interval still be narrower than the Lesson 27 z
      interval?
    \item What happens to $s$ when one fill is replaced by $560$ ml?
  \end{enumerate}

  \vspace{0.6em}
  \begin{block}{Connection to Lesson 29}
    The session can continue directly into a confidence interval for a
    population proportion and the $np$, $n(1-p)$ conditions.
  \end{block}

  \vfill
  \scriptsize Source: OpenStax, \textit{Introductory Business Statistics 2e},
  Chapter 8, Section 8.2. CC BY-NC-SA.
\end{frame}

\appendix



\end{document}
